\odiseachapter{Undécima parte}{Sirenas}

Poco cuesta imaginar un concilio de los olímpicos, similar a los que nos presenta Homero, en el que los inmortales debaten lo que hacer con Nolan.

---Al Tártaro con él ---exclama Atenea---. Es imperdonable que me haya convertido en una psiquiatra \emph{new age}, vestida de blanco.

---Deberíamos ponerlo de ayudante de Sísifo ---interviene Apolo---. A ver si aprende a no pintar a los héroes con brocha gorda.

---Con vuestro permiso, yo lo convertiría en algún animalito para hacerme compañía ---propone Circe---. A ver si se entera de cómo se debe tratar a las deidades.

---No seáis energúmenos ---dice Calipso---. Yo me lo quedo una temporada. El chaval no es feo ni tonto, y me sobra un poco de flor de loto.

---¿Y si le preguntamos a Ulises? ---propone Zeus.

Todos están de acuerdo y así lo mandan llamar. Huelga decir que Ulises está en el Olimpo, no en el Hades, donde su lugar lo ocupa el último modelo de \gr{εἴδωλον}, animado por IA. Zeus, muy amigo de formalidades, le dice:

\begin{verse}
Sangre Laercia, Odiseo milmañas, hijo del cielo,\\
dinos qué hacer con Christopher muypeliculero.\\
¿Lo salvarías tú o, por el contrario, implacable\\
y justiciero, lo condenarías al Averno?
\end{verse}

Odiseo se lo piensa un poco. Le tiene algo de tirria al esmerado director, por convertir un personaje tan complejo como él en un veterano de guerra con síndrome postraumático. Reconoce, por otra parte, que, siendo hijo del divinal Hollywood y de la ninfa Mercado, Nolan bienpopular tenía que ingeniárselas para pasar entre la Escila de las ventas en taquilla y la Caribdis de unos tiempos simplones, donde todo mensaje complejo se pierde en el vinoso ponto. Y lo cierto, se dice, es que la peli tampoco está tan mal. Si se le perdona:

\begin{verse}
Que Matt elprota pasmado se pase la aventura,\\
y al niño Telémaco no le dejen dar pie con bola,\\
que Penélope se crea Julieta y el machialtivo\\
pretendiente sea un vampiro disfrazado.

Peor aún, la bella Nausícaa, la hospitalaria,\\
y los mágicos feacios, su pueblo entero,\\
eliminado los ha, el muycicatero,\\
Polifemo, del ardiente ojo es tonto,\\
y al Pelida divino no ha dado aliento.
\end{verse}

La lista de ofensas es larga y Odiseo teme que si se pone a pensar le salga un poema épico de calamidades. Pero entonces se acuerda de una escena concreta y sentencia:

---Lo mandamos al Elíseo.

¿Qué escena es esa? El encuentro con los lestrigones, ya lo hemos dicho, está bien resuelto, si pasamos por alto la estética futurista de las armaduras de los gigantones ---de armaduras, dicho sea de paso, mejor no hablamos---, el paso entre los dos monstruos se salva con solvencia y también el episodio de las vacas del Sol. El viaje al Hades, si no tenemos en cuenta la carrera final ---los pobres aqueos se pasan la película corriendo--- y le perdonamos que se olvide de Aquiles, es de sobresaliente. Pero nada de eso le habría salvado el pellejo, dadas sus muchas ofensas.

Pero Ulises ha recordado a las sirenas, sabe cómo se sintió atado al mástil y está convencido de que Nolan ha sido capaz de capturar exactamente aquella sensación. Mejor aún, se da cuenta de que, en este caso, el cine es capaz de superar la palabra escrita, al menos durante unos instantes.

Hablando de palabra escrita, retomemos el viaje de nuestro héroe en el canto XII, cuando él y sus hombres regresan a la isla de Circe, después de haber escapado del Hades.

\begin{verse}
Mas cuando ya de aquel río Oceano dejó los raudales\\
nuestro navío y del mar anchimundo ganó el oleaje\\
hasta la isla de Eea, do casa y ruedo de baile\\
tiene la Aurora temprana y se hace del Sol el levante,\\
a ella llegados, hicimos que el barco en la playa encallase\\
y descendimos sobre el rompeolas a los arenales;\\
allí durmiendo aguardamos a que la divina alborease.
\end{verse}

Lo primero que hace Ulises es ir a buscar el cuerpo de Elpénor y rendirle honras fúnebres. Circe aparece enseguida y les dice:

\begin{verse}
«Qué pobres, sí, que con vida bajasteis a casa de Hades:\\
dos veces muertos, cuando una vez mueren los otros mortales.\\
Venga, comida comed y del vino bebed lo que falte\\
ya de este día, que, apenas el alba amanezca, la nave al\\
mar volveréis, y os diré yo el camino y de cada detalle\\
seña os daré, porque no en tal o cual mala urdimbre de males\\
os enredéis y en el mar o en la tierra se os venga desastre».
\end{verse}

Los aqueos no se hacen de rogar ---están muertos de hambre, nunca mejor dicho, después del susto del Hades--- y, tras cenar a sus anchas, se echan a dormir. Circe aprovecha para hablar con Odiseo.

\begin{verse}
Tomó ella entonces mi mano y, sentándonos lejos de aquellos,\\
a hacer preguntas se puso, y todo quería saberlo;\\
y de eso todo, a verdad, por menudo yo le hice recuento.\\
Y entonces ya ánima Circe hablándome estaba de nuevo:

Salvado está ---y bien salvado--- todo eso; escúchame ahora\\
lo que te voy a decir, y que un dios dejará en tu memoria.\\
Con las Sirenas darás lo primero, que hechizan y emboban\\
a cuantos hombres ---da igual quiénes sean--- con ellas se topan.\\
Que a todo el que en su ignorancia dé allí y la voz blanca oiga\\
de las Sirenas, ya más ni su prole ni ya más su esposa\\
a recibirlo saldrá, porque a casa él volvió, jubilosa,\\
que lo atraerán las Sirenas, de aquel su cantar silbadoras,\\
a esa pradera en que están y que tiene por playa a redonda\\
los huesos y los pellejos resecos de carnes hediondas.\\
Remad, remad, y el oído de todos tus hombres tapona\\
con dulce cera que habrás ablandado, y nadie allí oiga.

Y si es que a ti el corazón a escucharlas te llama, que pongan\\
cuerda alredor de tus manos y pies que al palo te coja\\
de la goleta bien tiesto, curando que no quede floja:\\
de las Sirenas la voz a placer así oirás, de una y otra.\\
Y si a tus hombres suplicas y mandas que suelten las sogas,\\
que echen más nudos y más cada vez en aquellas maromas.
\end{verse}

Después de dejar atrás a las sirenas, le esperan Escila y Caribdis. Pagado el mortal peaje, le espera la isla de Trinacia, donde campean las vacas del Sol, que ni él ni su tripulación deben tocar. También sabemos lo que le dice:

\begin{verse}
Si sin tocar se las dejas y atiendes no más al regreso,\\
a Ítaca entonces podríais llegar, aunque penas sufriendo;\\
pero si se las tocáis, desde aquí ya te auguro hundimiento\\
para tu nave y tus hombres. Y aun cuando salieras tú ileso,\\
tarde, y no bien, llegarías, y todos tus hombres ya muertos.
\end{verse}

Dicho todo lo cual, la beltrenza se despide, mandándoles un buen viento. ¡Oh, Circe! Después de un largo año juntos, después de mandarle a visitar los infiernos, ¿dejaste ir, sin más, a Odiseo? ¿O, una vez más, el mientisagaz nos está ocultando algo?

Ya veremos más adelante. Por ahora, la nave zarpa y Ulises pone a sus hombres sobre aviso:

\begin{verse}
«Bien no está, amigos, que solo uno o dos de nosotros conozcan\\
esas anuncias que Circe me dio, la divina entre diosas;\\
os las diré, y que sabiéndolas todos nos venga, si toca,\\
la muerte, o de ella salgamos hurtándonos a nuestra sombra.\\
De las Sirenas malhechiceras la voz sibilosa\\
es lo primero que manda esquivar y la braña en que moran.\\
Que solo yo las escuche, me dijo, y para eso una soga\\
difícil anudaréis, que allí de firme me coja\\
al pie del mástil bien prieto, mirando que no quede floja.\\
Y si os suplico y ordeno que me desatéis de la soga,\\
lazada y nueva lazada por cima echaréis de las otras».
\end{verse}

Y así llega la negra nave a la isla maldita:

\begin{verse}
Y de repente el viento cesó y un silencio muy quieto\\
se hizo en el aire encalmado, y el mar adormíalo el cielo.\\
Puestos en pie, recogieron mis hombres la vela y la dieron\\
a la cubierta del cuenco navío, y, sentándose al remo,\\
el agua blanca volvían del ir de los lisos maderos.\\
Yo hundí mi espada en el gran pan de cera y me lo iba partiendo\\
en migas que apretujaba entre mis manos de recio;\\
y entró enseguida la cera en calor, la obligaba mi empeño\\
y el soberano Hiperiónida, el Sol, que mandaba su fuego:\\
uno por uno tapé los oídos a mis compañeros.

Ellos me ataron de manos y pies en la nave, bien prieto\\
al bajo mismo del palo, y al palo otros nudos le hicieron;\\
sentados luego a los remos, batían ya el mar caniciento.\\
Como a una voz, si uno grita, andaríamos de esas de lejos\\
en nuestro vuelo, cuando ellas notaron que había un velero\\
margalopando allí cerca, y se abrió el sibilante coreo:
\end{verse}

¿Qué le dicen las Sirenas? Homero escribe:

\begin{verse}
«Ven, alabado, aquí ven, de tu pueblo, Odiseo, alta gloria,\\
trae esa nave y la voz óyenos de la una y la otra.\\
Que hombre ninguno en su negra goleta se fue de nosotras\\
sin escuchar esta voz melidulce de nuestras bocas,\\
que la gozó, y prosiguió su viaje, más sabio de cosas:\\
todo sabemos las dos del penar que en la muy ancha Troya\\
llevaron teucros y argivos por gana de dioses y diosas,\\
todo las dos de la tierra sabemos de bienes dadora».
\end{verse}

\begin{verse}
Así decían, dejando ir hermosa su voz; en mi pecho\\
solo quería ya oír, y movía las cejas pidiendo\\
que me soltaran; y más a los remos postrábanse aquellos.\\
Y se pusieron en pie Perimedes y Euríloco y fueron\\
a echar más nudos por sobre los nudos, y más me prendieron.\\
Cuando por fin las dejamos atrás y ya nada de eso\\
---canto ni voz de Sirenas--- oíamos, mis compañeros,\\
mis fieles, se despegaban la cera que unté bien adentro\\
de sus oídos, y a mí de mis lazos librábanme prestos.
\end{verse}

Y no hay más. Por una vez, el Poeta se limita a contarnos lo que pasa y no nos muestra cómo se siente Odiseo, o en qué consiste la magia de las Sirenas.

Esa era otra de las cuentas que yo tenía pendientes con el granbardo. En \emph{Abandonando Ítaca}\footnote{\url{https://eolasediciones.es/libro/abandonando-itaca-leaving-ithaca/}}, escribo mi propia versión. Homero no da muchos detalles de cómo se sentían los protagonistas de su historia. Los muestra a menudo comiendo o bebiendo ---generalmente en exceso---, discutiendo, llorando, vociferando, o intentando salvar el pellejo. Pero rara vez sentimos sus emociones, rara vez llegamos a ponernos en su piel. Y sin embargo, son hombres como nosotros:

\begin{verse}
Ese sentido de lo sagrado. Tú y tu banda de bárbaros\\
lo teníais. A menudo, de noche, cuando la luna estaba alta y el viento calmo,\\
los hombres se callaban. Al cabo de un rato, se alzaba una voz,\\
cantando una canción triste.\\
Podría ser por la muerte de Héctor, o la muerte de Aquiles,\\
podía recordar las naves perdidas o invocar las llamas que incendiaban Troya,\\
lamentar el destino de Patroclo o añorar la belleza de Helena,\\
pero siempre era triste. Los hombres están en su mejor momento,\\
cuando cantan sobre lo que han perdido.
\end{verse}

Y es, justamente, de ese sentido de lo sagrado de lo que se alimentan las sirenas:

\begin{verse}
Quizás esto es lo que hizo que las sirenas\\
fuesen tan temibles y tan peligrosas,\\
cantaban simplemente canciones de los hombres\\
con acordes más puros.\\
Oh, cómo lloraste atado al mástil de la embarcación,\\
escuchando su voz, mientras tu sorda tripulación se dedicaba a sus asuntos,\\
ignorando su hechizo. Y, sin embargo, lo volverías a hacer mil veces más.\\
Si tu viaje tenía alguna justificación,\\
si pasar entre Escila y Caribdis valía el riesgo,\\
si luchar contra los cíclopes y los lestrigones tenía sentido,\\
era tan solo para darte la oportunidad\\
de escuchar esas voces,\\
recordando todo lo que en la vida merece la pena.
\end{verse}

¿Y qué valía la pena? En sus años de vejez, Ulises no se acuerda de la gloria de la batalla, ni de las joyas robadas, ni de los palacios y los banquetes, sino de cosas mucho más sencillas. Son esas cosas las que cantan las sirenas. Son esas cosas las que parten el alma del marino y le hacen desear arrojarse a las olas:

\begin{verse}
Aquellos primeros años, cuando la risa de Penélope\\
llenaba la casa de alegría,\\
con el bebé, durmiendo entre tus fuertes brazos,\\
sus pequeñas manos en torno a tu cuello, más preciosas\\
que un collar de perlas corintias.\\
El amanecer pintando el cielo de carmesí,\\
tus pulmones llenándose con el aire fresco de un nuevo día.\\
El sabor de las aceitunas en tu boca, el placer del pan y el vino\\
tras un día de duro trabajo. Cantaron todo eso, y más,\\
cantaron ríos, y montañas nevadas que nunca has visto,\\
todas vacías de dioses, pero llenas de gracia.
\end{verse}

Vacías de dioses, pero llenas de gracia. Odiseo ha necesitado llegar a la senectud para entender el canto de las sirenas. Y Nolan consigue captarlo en unas escenas bellísimas, en las que un hombre atado a un mástil grita desesperado, mientras la música de las esferas, arañándole el espíritu, le recuerda lo que ha perdido:

\begin{verse}
Cantaron todos esos incontables momentos\\
de simple felicidad.\\
¿Adónde fueron todas esas mañanas?\\
¿Qué fue de aquella joven pareja que caminaba, cogida de la mano,\\
por las mismas colinas que ahora recorres en soledad?\\
Sí, las sirenas cantaban sobre lo sagrado,\\
lo sagrado que no necesita dioses ni sacrificios,\\
lo sagrado, que se teje, como el sudario de Penélope,\\
sobre la vida de todos los hombres. Mañanas, y ríos, y niños, besos,\\
y estrellas y aceitunas,\\
pan y vino, y tiempo, tiempo que fluye,\\
huyendo.
\end{verse}
